% !PRINTED 438
\weberpage{438}
%<<CONT
\[
h = 2m,\ 4m,\ 6m,\ \ldots (n-1)m
\]
annimmt, so erhält man in (21) weder zwei gleiche, noch zwei entgegengesetzte Werthe, und es ergeben sich alle von einander verschiedenen Werthe von $x^2$.

Es wird also für ein \gesperrt{ungerades} $n$
\[
B_n(x) = \prod\limits_{1,\,\frac{n-1}{2}}^{\nu} \left(x^2 - 4\cos^2\frac{2\nu m\pi}{n}\right). \tag{24}
\]

Setzen wir in dieser Formel
\[
x = 2\cos\varphi,\quad x^2 = 4(1-\sin^2\varphi),\quad \cos^2\frac{2\nu m\pi}{n} = 1 - \sin^2\frac{2\nu m\pi}{n},
\]
so können wir sie so darstellen:
\[
\frac{\sin n\varphi}{\sin\varphi} = (-1)^{\frac{n-1}{2}} 2^{n-1} \prod\limits_{1,\,\frac{n-1}{2}}^{\nu} \left(\sin^2\varphi - \sin^2\frac{2\nu m\pi}{n}\right), \tag{25}
\]
woraus wieder für $\varphi=0$ die Formel (20) folgt.

Da $\cos n\varphi$ bei ungeradem $n$ verschwindet, wenn
\[
\cos\varphi = \pm \sin\frac{2\nu m\pi}{n}
\]
ist, so haben wir damit auch für diesen Fall die Factorenzerlegung der Function $A_n(x)$ gefunden. Es ist für \gesperrt{ungerades} $n$
\[
A_n(x) = x \prod\limits_{1,\,\frac{n-1}{2}}^{\nu} \left(x^2 - 4\sin^2\frac{2\nu m\pi}{n}\right). \tag{26}
\]

Die Grössen
\[
\alpha = \sin^2\frac{2\nu m\pi}{n} \tag{27}
\]
sind also die Wurzeln einer algebraischen Gleichung vom Grade $(n-1):2$
\[
\Phi_n(x) = 0,
\]
und wir erhalten die Function $\Phi_n(x)$, wenn wir an Stelle der Variablen $x^2$ in der Function $A_n(x):x$ eine neue Variable $4x$ setzen und durch $2^{n-1}$ dividiren. Es wird daher [Formel (11)]
\[
\Phi_n(x) = \overset{\alpha}{\Pi}(x-\alpha)
= x^{\frac{n-1}{2}} - \frac{n}{4} x^{\frac{n-3}{2}} + \frac{n(n-3)}{1.2}\,\frac{1}{16}\, x^{\frac{n-5}{2}} - \cdots, \tag{28}
\]
und für (25) können wir schreiben
\[
\frac{\sin n\varphi}{\sin\varphi}
= (-1)^{\frac{n-1}{2}} 2^{n-1} \overset{\alpha}{\Pi}(\sin^2\varphi - \alpha). \tag{29}
\]
